\documentclass[12pt]{article}
\usepackage[utf8]{inputenc}
\usepackage[T1]{fontenc}
\usepackage[margin=1in]{geometry}
\usepackage{setspace}
\usepackage{xcolor}
\usepackage[colorlinks=true, urlcolor=blue!60!black, linkcolor=blue!60!black]{hyperref}
\setlength{\parindent}{2em}
\setlength{\parskip}{1.25em}
\renewcommand{\baselinestretch}{1.0}
\begin{document}
\begin{center}
{\LARGE \textit{Hans v. Louisiana}, 134 U.S. 1 (1890)\par}
\vspace{0.8em}
{\large Case Brief\par}
\end{center}
\vspace{1em}

\section{Memory Jogger}
Despite the Eleventh Amendment's text referencing only suits by out-of-state or foreign citizens, the background constitutional principle of state sovereign immunity bars a citizen from suing his own state in federal court without its consent, even on a federal-question claim.

\section{Detailed Case Facts}
Under an act of the Louisiana legislature approved on January 24, 1874, the State of Louisiana issued ``consolidated bonds'' payable forty years from January 1, 1874, with semiannual interest represented by attached coupons (\href{https://supreme.justia.com/cases/federal/us/134/1/}{134 U.S. at 1--2}; \href{https://www.courtlistener.com/opinion/92702/hans-v-louisiana/}{CourtListener}). By a companion constitutional amendment adopted in 1874, Louisiana declared those consolidated bonds to be a valid, unimpeachable contract between the State and every bondholder and dedicated an annual tax to pay principal and interest without further legislative appropriation (\href{https://supreme.justia.com/cases/federal/us/134/1/}{134 U.S. at 2}).

Following the end of Reconstruction, Louisiana adopted a new state constitution in 1879 that remitted the consolidated bond coupons falling due on January 1, 1880, transferred the interest taxes already collected for those coupons to defray general state government expenses, and superseded the 1868 Constitution and its 1874 amendments (\href{https://supreme.justia.com/cases/federal/us/134/1/}{134 U.S. at 2--3}; \href{https://texaslawreview.org/state-sovereign-immunity-after-the-revolution/}{Texas Law Review}). Bernard Hans, a citizen of Louisiana who held consolidated bond coupons maturing on January 1, 1880, filed suit against the State of Louisiana in December 1884 seeking \$87,500 plus legal interest from January 1, 1880, and costs (\href{https://supreme.justia.com/cases/federal/us/134/1/}{134 U.S. at 1, 3}). Hans alleged that Louisiana's 1879 constitutional repudiation impaired the obligation of his bond contract in violation of the Contract Clause, Article I, Section 10 of the United States Constitution, and invoked federal-question jurisdiction under Article III, Section 2 and the Judiciary Act of March 3, 1875 (\href{https://supreme.justia.com/cases/federal/us/134/1/}{134 U.S. at 3, 9}).

\section{Procedural History}
Hans commenced his action in the United States Circuit Court for the Eastern District of Louisiana in December 1884, serving a citation on the Governor of Louisiana (\href{https://supreme.justia.com/cases/federal/us/134/1/}{134 U.S. at 1, 3}). The Attorney General of Louisiana filed an exception asserting that the court lacked jurisdiction \textit{ratione personae} because the State could not be sued without its consent (\href{https://supreme.justia.com/cases/federal/us/134/1/}{134 U.S. at 3}). In 1885, the Circuit Court (opinion by Judge Edward Coke Billings) sustained the State's exception and dismissed the suit, reported at \textit{Hans v. Louisiana}, 24 F.\ 55 (C.C.E.D.\ La.\ 1885) (\href{https://supreme.justia.com/cases/federal/us/134/1/}{134 U.S. at 4}; \href{https://texaslawreview.org/state-sovereign-immunity-after-the-revolution/}{Texas Law Review}). Hans brought a writ of error to the United States Supreme Court (\href{https://supreme.justia.com/cases/federal/us/134/1/}{134 U.S. at 1, 9}). The case was argued and submitted on January 22, 1890, and decided on March 3, 1890, with the Supreme Court unanimously affirming the dismissal (\href{https://supreme.justia.com/cases/federal/us/134/1/}{134 U.S. at 1, 21}).

\section{Judicial Votes}
\textbf{Vote:} 9--0 (8--1 on reasoning), affirming the Circuit Court's dismissal for lack of jurisdiction (\href{https://supreme.justia.com/cases/federal/us/134/1/}{134 U.S. at 9, 21}).
\begin{itemize}
\item \textbf{Majority (8):} Justice Joseph P. Bradley (author), joined by Chief Justice Melville W. Fuller, Justice Samuel F. Miller, Justice Stephen J. Field, Justice Horace Gray, Justice Samuel Blatchford, Justice Lucius Q. C. Lamar, and Justice David J. Brewer (\href{https://supreme.justia.com/cases/federal/us/134/1/}{Justia}).
\item \textbf{Concurrence (1):} Justice John Marshall Harlan (author), concurring in the judgment only (\href{https://supreme.justia.com/cases/federal/us/134/1/}{134 U.S. at 21}).
\item \textbf{Dissents (0):} None.
\end{itemize}

\section{Holding}
\textbf{Judgment:} Affirmed. The Circuit Court properly dismissed Bernard Hans's suit against the State of Louisiana (\href{https://supreme.justia.com/cases/federal/us/134/1/}{134 U.S. at 21}).

\textbf{Rule Announced:} An unconsenting State cannot be sued in a federal court by one of its own citizens, even when the suit arises under the Constitution or laws of the United States (\href{https://supreme.justia.com/cases/federal/us/134/1/}{134 U.S. at 1, 10--15, 21}). Neither Article III's federal-question grant nor the Judiciary Act of 1875 authorizes private suits against an unconsenting State, because the Constitution presupposes the preexisting sovereign immunity of the States as an attribute of sovereignty not surrendered in the plan of the Convention (\href{https://supreme.justia.com/cases/federal/us/134/1/}{134 U.S. at 13--19}).

\textbf{Holding vs.\ Dicta:} The core holding is that a citizen cannot sue his own non-consenting State in a federal Circuit Court on a federal-question claim (here, a Contract Clause claim on state bonds). Two major passages in Justice Bradley's opinion are dicta: first, his sweeping disapproval of the 1793 majority in \textit{Chisholm v. Georgia}, 2 U.S.\ (2 Dall.) 419 (1793), and endorsement of Justice Iredell's dissent---which Justice Harlan specifically flagged as unnecessary to the decision (\href{https://supreme.justia.com/cases/federal/us/134/1/}{134 U.S. at 12--14, 21}); and second, Bradley's reassuring observation that while a State cannot be compelled by suit to perform its contracts, ``any attempt on its part to violate property or rights acquired under its contracts may be judicially resisted'' (\href{https://supreme.justia.com/cases/federal/us/134/1/}{134 U.S. at 20--21}).

\section{Analysis of Opinions}

\subsection{Majority Opinion (Justice Bradley)}
Justice Bradley began by acknowledging Hans's textual argument: the Eleventh Amendment literally bars only suits ``commenced or prosecuted against one of the United States by Citizens of another State, or by Citizens or Subjects of any Foreign State,'' and says nothing about suits by a State's own citizens under Article III's federal-question clause (\href{https://supreme.justia.com/cases/federal/us/134/1/}{134 U.S. at 10}). Yet Bradley rejected a literal reading as producing an ``anomalous result''---allowing a State to be sued in federal court by its own citizens on federal claims even though it could not be sued on identical claims by out-of-state or foreign citizens (under \textit{Louisiana v. Jumel}, \textit{Hagood v. Southern}, and \textit{In re Ayers}) and did not permit itself to be sued in its own courts (\href{https://supreme.justia.com/cases/federal/us/134/1/}{134 U.S. at 10--11}).

To justify departing from the bare text, Bradley turned to Founding-era history and the aftermath of \textit{Chisholm v. Georgia}, 2 U.S.\ (2 Dall.) 419 (1793). \textit{Chisholm}, he wrote, created ``such a shock of surprise throughout the country'' that Congress and the States swiftly adopted the Eleventh Amendment, which did not merely create a new textual exception but declared that the Constitution ``shall not be construed'' to authorize private suits against States---effectively vindicating Justice Iredell's lone dissent in \textit{Chisholm} (\href{https://supreme.justia.com/cases/federal/us/134/1/}{134 U.S. at 11--12}). Bradley quoted Alexander Hamilton in \textit{Federalist No.\ 81} (``It is inherent in the nature of sovereignty not to be amenable to the suit of an individual without its consent'') alongside assurances by James Madison and John Marshall at the Virginia Ratifying Convention that individuals could not hail sovereign States into federal court (\href{https://supreme.justia.com/cases/federal/us/134/1/}{134 U.S. at 12--14}). Asking whether the States would ever have ratified the Eleventh Amendment had it included a proviso allowing a State's own citizens to sue it in federal court, Bradley declared the suggestion ``almost an absurdity on its face'' (\href{https://supreme.justia.com/cases/federal/us/134/1/}{134 U.S. at 15}).

Bradley reinforced this constitutional presumption with three additional points. First, citing \textit{Beers v. Arkansas}, \textit{Briscoe v. Bank of Kentucky}, and \textit{Cunningham v. Macon \& Brunswick R.R.}, he reiterated that ``[t]he suability of a State without its consent was a thing unknown to the law'' (\href{https://supreme.justia.com/cases/federal/us/134/1/}{134 U.S. at 16--18}). Second, as a statutory ground, he noted that the Judiciary Act of 1875 gave Circuit Courts federal-question jurisdiction ``concurrent with the courts of the several States''; because state courts could not entertain unconsented suits against the State, Congress could not have intended to grant federal courts a novel jurisdiction unknown to state courts (\href{https://supreme.justia.com/cases/federal/us/134/1/}{134 U.S. at 18--19}). Third, he brushed aside Chief Justice Marshall's contrary suggestion in \textit{Cohens v. Virginia}, 19 U.S.\ (6 Wheat.) 264, 412 (1821), as ``unnecessary to the decision, and in that sense extrajudicial'' (\href{https://supreme.justia.com/cases/federal/us/134/1/}{134 U.S. at 19--20}). While noting that a State's contracts rest on its political honor and that active invasions of property rights under state contracts may still be resisted defensively, Bradley held that courts may not compel a State to pay its public debts (\href{https://supreme.justia.com/cases/federal/us/134/1/}{134 U.S. at 20--21}).

\subsection{Concurring Opinion (Justice Harlan)}
Justice Harlan wrote a brief, pointed concurrence agreeing solely with the bottom-line holding that federal judicial power does not extend to a direct suit against a State by one of its own citizens without the State's consent (\href{https://supreme.justia.com/cases/federal/us/134/1/}{134 U.S. at 21}). However, Harlan refused to join ``many things said in the opinion,'' specifically objecting to Bradley's criticism of \textit{Chisholm v. Georgia} (\href{https://supreme.justia.com/cases/federal/us/134/1/}{134 U.S. at 21}). In Harlan's view, Bradley's attack on \textit{Chisholm} was both unnecessary to decide \textit{Hans} and wrong on the merits, because \textit{Chisholm} ``was based upon a sound interpretation of the Constitution as that instrument then was'' (\href{https://supreme.justia.com/cases/federal/us/134/1/}{134 U.S. at 21}).

\subsection{Dissenting Opinions}
No dissenting opinions were filed.

\section{Examples: Future Applications}
Applying the sovereign-immunity principle of \textit{Hans v. Louisiana} (and its doctrinal boundaries) produces the following five applications:

\begin{enumerate}
\item \textbf{Same-Side Example 1 (Barred Under \textit{Hans}):} A resident of Texas sues the State of Texas directly in federal district court for retrospective money damages, alleging that a state statute violates her federal statutory overtime rights enacted under Congress's Article I Commerce Clause power. \textit{Why:} Just as in \textit{Hans}, a citizen is suing her own unconsenting State in federal court on a federal-question claim. Under \textit{Hans}---later reinforced by \href{https://texaslawreview.org/state-sovereign-immunity-after-the-revolution/}{\textit{Seminole Tribe of Florida v. Florida}, 517 U.S. 44 (1996)}---sovereign immunity bars the suit, and Congress may not abrogate that immunity using its Article I commerce power.

\item \textbf{Same-Side Example 2 (Barred Under \textit{Hans}):} A Louisiana state employee sues the Louisiana Department of Revenue (an arm of the State) in federal court seeking back pay for an alleged breach of a state pension statute, claiming a federal Contract Clause violation. \textit{Why:} The claim mirrors \textit{Hans} directly: an in-state citizen seeks a coercive monetary judgment payable out of the state treasury against an unconsenting State on a Contract Clause theory (\href{https://supreme.justia.com/cases/federal/us/134/1/}{134 U.S. at 3, 20--21}; \href{https://constitution.congress.gov/browse/essay/amdt11-6-3/ALDE_00013684/}{Constitution Annotated}).

\item \textbf{Opposite-Side Example 1 (Permitted Under \textit{Hans}'s Boundaries):} A Louisiana citizen sues the Louisiana Attorney General in his official capacity in federal court seeking only a prospective injunction to stop him from enforcing an unconstitutional state law. \textit{Why:} \textit{Hans} bars direct suits against the State itself (and suits seeking retrospective funds from the state treasury), but it expressly distinguished suits against individual officers to resist unlawful invasions of rights (\href{https://supreme.justia.com/cases/federal/us/134/1/}{134 U.S. at 16, 20--21}). Under the officer-suit doctrine later crystallized in \textit{Ex parte Young}, 209 U.S. 123 (1908), prospective injunctive relief against a state official to halt an ongoing violation of federal law is not barred (\href{https://constitution.heritage.org/essays/amdt-11/}{Heritage Guide}; \href{https://www.law.cornell.edu/constitution-conan/amendment-11/exceptions-to-eleventh-amendment-immunity-officer-suits}{Cornell LII}).

\item \textbf{Opposite-Side Example 2 (Permitted Under \textit{Hans}'s Boundaries):} A citizen of Ohio sues a municipal county (such as Franklin County) in federal court for defaulting on county-issued municipal bonds, or the United States itself sues the State of Ohio in federal court to enforce a federal statute. \textit{Why:} State sovereign immunity under \textit{Hans} protects the State itself and its arms, not political subdivisions such as counties or municipalities (\textit{Lincoln County v. Luning}, 133 U.S. 529 (1890), decided the very same day as \textit{Hans}), and States surrendered immunity to suits brought by the United States in the plan of the Convention (\href{https://texaslawreview.org/state-sovereign-immunity-after-the-revolution/}{Texas Law Review}).

\item \textbf{Fence-Sitter Example (Borderline Under \textit{Hans}):} Congress enacts a federal statute under its Article I war powers or eminent-domain power (or delegates federal condemnation authority to a private pipeline company) expressly authorizing private citizens to sue an unconsenting State in federal or state court. \textit{Why:} On the \textit{Hans} side, a private citizen is hailing an unconsenting State into court under an Article I statute, which \textit{Hans} and \textit{Seminole Tribe} ordinarily forbid. On the opposite side, Hamilton's \textit{Federalist No.\ 81} passage quoted in \textit{Hans} expressly carves out immunity surrendered ``in the plan of the convention'' (\href{https://supreme.justia.com/cases/federal/us/134/1/}{134 U.S. at 13})---an exception the modern Court has narrowly invoked in divided 5--4 decisions for federal bankruptcy (\textit{Central Virginia Community College v. Katz}), delegated federal eminent domain (\textit{PennEast Pipeline Co. v. New Jersey}), and war powers (\textit{Torres v. Texas Department of Public Safety}) (\href{https://texaslawreview.org/state-sovereign-immunity-after-the-revolution/}{Texas Law Review}).
\end{enumerate}

\section{Critique}
\textbf{Originalist and Textualist Perspectives:} \textit{Hans v. Louisiana} has provoked one of the fiercest textualist and originalist debates in American constitutional law. Strict textualists and ``diversity theory'' scholars---such as William Fletcher, John Gibbons, and (in a modified form) Akhil Amar---argue that \textit{Hans} flatly contradicts the text of both Article III and the Eleventh Amendment: the Eleventh Amendment was drafted narrowly to repeal only \textit{Chisholm}'s Citizen-State Diversity Clause jurisdiction, leaving federal-question jurisdiction over suits by a State's own citizens untouched (\href{https://texaslawreview.org/state-sovereign-immunity-after-the-revolution/}{Texas Law Review}; \href{https://ndlawreview.org/wp-content/uploads/2013/07/Menashi.pdf}{Notre Dame Law Review}). Conversely, modern ``original-law'' originalists and structural textualists---notably William Baude, Stephen E. Sachs, and Caleb Nelson---defend \textit{Hans}'s result on different grounds: they contend that sovereign immunity was a background doctrine of general or common law (analogous to personal jurisdiction) that Article III never abrogated and that the Eleventh Amendment merely prevented courts from misconstruing Article III to abolish it (\href{https://lawreview.gmu.edu/print__issues/sovereign-immunity-and-the-two-tiers-of-article-iii/}{George Mason Law Review}; \href{https://texaslawreview.org/state-sovereign-immunity-after-the-revolution/}{Texas Law Review}). Even among these defenders, however, scholars note that if sovereign immunity is a background common-law rule rather than a rigid constitutional bar, Congress might possess greater power to abrogate it than the Rehnquist Court later allowed in \textit{Seminole Tribe} and \textit{Alden v. Maine} (\href{https://texaslawreview.org/state-sovereign-immunity-after-the-revolution/}{Texas Law Review}; \href{https://constitution.heritage.org/essays/amdt-11/}{Heritage Guide}).

\textbf{Progressive and Doctrinal Perspectives:} Progressive and federal-courts scholars---including Vicki Jackson, Edward A. Purcell, Jr., and John V. Orth---critique \textit{Hans} on both historical-realist and rule-of-law grounds (\href{https://texaslawreview.org/state-sovereign-immunity-after-the-revolution/}{Texas Law Review}). Historically, these critics emphasize that \textit{Hans} was not a neutral restoration of Founding principles, but a pragmatic post-Reconstruction retreat: after the Compromise of 1877, white supremacist ``Redeemer'' governments across the South repudiated Reconstruction-era bonds, and the Supreme Court lacked the political will or enforcement capacity to compel Southern legislatures to levy taxes to pay bondholders (\href{https://texaslawreview.org/state-sovereign-immunity-after-the-revolution/}{Texas Law Review}; \href{https://www.law.cornell.edu/constitution-conan/amendment-11/postbellum-jurisprudence-on-the-eleventh-amendment}{Cornell LII}). Doctrinally, critics argue that elevating \textit{Hans} into an unassailable constitutional shield---as the Court later did in \textit{Welch v. Texas Department of Highways} (\href{https://supreme.justia.com/cases/federal/us/483/468/}{483 U.S. 468 (1987)}), \textit{Seminole Tribe}, and \textit{Alden}---undermines the Supremacy Clause and the remedial maxim of \textit{Marbury v. Madison} by leaving citizens without a damages remedy when States violate federal constitutional or statutory rights (\href{https://www.law.cornell.edu/supct/html/94-12.ZD.html}{Cornell LII}; \href{https://texaslawreview.org/state-sovereign-immunity-after-the-revolution/}{Texas Law Review}).

\textbf{Underlying Competing Values:} \textit{Hans v. Louisiana} pits \textit{state fiscal autonomy and sovereign dignity} (shielding state treasuries and democratic budget priorities from coercive judicial judgments) against \textit{constitutional supremacy and individual accountability} (ensuring that States cannot violate the U.S.\ Constitution or federal law with financial impunity).

\section{Key Quotations}
\begin{quote}
``Can we suppose that, when the Eleventh Amendment was adopted, it was understood to be left open for citizens of a State to sue their own state in the federal courts, whilst the idea of suits by citizens of other states, or of foreign states, was indignantly repelled? Suppose that Congress, when proposing the Eleventh Amendment, had appended to it a proviso that nothing therein contained should prevent a State from being sued by its own citizens in cases arising under the Constitution or laws of the United States, can we imagine that it would have been adopted by the States? The supposition that it would is almost an absurdity on its face.''
\end{quote}
\noindent --- Justice Joseph P. Bradley, Opinion of the Court, \href{https://supreme.justia.com/cases/federal/us/134/1/}{134 U.S. at 15} (\href{https://www.courtlistener.com/opinion/92702/hans-v-louisiana/}{CourtListener}).

\begin{quote}
``The suability of a State without its consent was a thing unknown to the law.''
\end{quote}
\noindent --- Justice Joseph P. Bradley, Opinion of the Court, \href{https://supreme.justia.com/cases/federal/us/134/1/}{134 U.S. at 16} (\href{https://www.courtlistener.com/opinion/92702/hans-v-louisiana/}{CourtListener}).

\begin{quote}
``It is inherent in the nature of sovereignty not to be amenable to the suit of an individual without its consent. This is the general sense and the general practice of mankind, and the exemption, as one of the attributes of sovereignty, is now enjoyed by the government of every State in the Union. Unless, therefore, there is a surrender of this immunity in the plan of the convention, it will remain with the States, and the danger intimated must be merely ideal.''
\end{quote}
\noindent --- Justice Joseph P. Bradley, Opinion of the Court (quoting Alexander Hamilton, \textit{The Federalist No.\ 81}), \href{https://supreme.justia.com/cases/federal/us/134/1/}{134 U.S. at 13} (\href{https://www.courtlistener.com/opinion/92702/hans-v-louisiana/}{CourtListener}).

\begin{quote}
``While the State cannot be compelled by suit to perform its contracts, any attempt on its part to violate property or rights acquired under its contracts may be judicially resisted, and any law impairing the obligation of contracts under which such property or rights are held is void and powerless to affect their enjoyment.''
\end{quote}
\noindent --- Justice Joseph P. Bradley, Opinion of the Court, \href{https://supreme.justia.com/cases/federal/us/134/1/}{134 U.S. at 20--21} (\href{https://www.courtlistener.com/opinion/92702/hans-v-louisiana/}{CourtListener}).

\begin{quote}
``I concur with the Court in holding that a suit directly against a State by one of its own citizens is not one to which the judicial power of the United States extends, unless the State itself consents to be sued. Upon this ground alone, I assent to the judgment. But I cannot give my assent to many things said in the opinion. The comments made upon the decision in Chisholm v. Georgia do not meet my approval. They are not necessary to the determination of the present case. Besides, I am of opinion that the decision in that case was based upon a sound interpretation of the Constitution as that instrument then was.''
\end{quote}
\noindent --- Justice John Marshall Harlan, Concurring, \href{https://supreme.justia.com/cases/federal/us/134/1/}{134 U.S. at 21} (\href{https://www.courtlistener.com/opinion/92702/hans-v-louisiana/}{CourtListener}).

\end{document}
